\documentclass[11pt]{article}
\usepackage[T1]{fontenc}
\usepackage[utf8]{inputenc}
\usepackage[margin=1in]{geometry}
\usepackage{amsmath,amssymb}
\usepackage{booktabs}
\usepackage{enumitem}
\usepackage{underscore}      % lets underscores print literally in text / \texttt
\usepackage[hidelinks]{hyperref}
\setlist[itemize]{leftmargin=1.4em}
\newcommand{\code}[1]{\texttt{#1}}
\newcommand{\Fval}[1]{{\small(\,$F\!\approx\!#1$\,)}}

\title{\textbf{Feature Documentation — Merged 2-Class Model (v16)}\\[2pt]
\large soft vs hard}
\author{WalkSensePlace}
\date{}

\begin{document}
\maketitle

Every feature is computed \textbf{per step} (one gait cycle, from one total-pressure peak to the
next, ${\sim}0.5$–$1.2$\,s at $f_s=62.5$\,Hz). Sensors used: 12 plantar pressure pads, a 3-axis
accelerometer, and a 3-axis gyroscope (the magnetometer is dead and unused). Steps that span a
surface change or a \code{transition}/\code{wait} tag are excluded, so each step lies on a single
surface.

Two conventions appear throughout:
\begin{itemize}
  \item \textbf{``norm''\,/\,session-normalised} --- divided by a within-recording scale (e.g.\ that
  step's mean total pressure) so the value doesn't scale with body weight or how tightly the insole
  was worn.
  \item \textbf{amplitude-invariant} --- the metric is a ratio or shape number, unaffected by
  overall size.
\end{itemize}

The target is \code{soft} (grass / compliant) vs \code{hard} (dirt, compact ground, concrete,
asphalt, exposed aggregate, brick). The model uses \textbf{48 features}; the $F$ values are ANOVA
soft-vs-hard separation scores (higher $=$ better). Person-confound features that are computed but
\textbf{dropped} are listed at the end.

\medskip\noindent\textbf{Why the soft/hard split is the reliable product.} These features primarily
encode surface \emph{compliance}. On the soft-vs-hard target the top features separate the classes
about twice as sharply as on the 3-class target --- e.g.\ the strongest,
\code{orient_rest_a_roll3_std}, scores $F\approx3109$ here vs $\approx1612$ in the 3-class model.
The features needed to split the \emph{hard} class further (firm natural ground vs pavement) are
weak, which is why the 3-class \code{loose} label doesn't generalise but soft-vs-hard does.

\section*{Notation}
For one step of length $L$ samples, let $p_t$ be total pressure; $s_{i,t}$ the value of pad $i$
($i=1,\dots,12$); heel sum $h_t=\sum_{i\in\mathcal H}s_{i,t}$ and forefoot sum
$f_t=\sum_{i\in\mathcal F}s_{i,t}$ over the heel set $\mathcal H$ and forefoot set $\mathcal F$.
Accelerometer axes are $a^x_t,a^y_t,a^z_t$ with magnitude
$\lVert a_t\rVert=\sqrt{(a^x_t)^2+(a^y_t)^2+(a^z_t)^2}$; gyroscope $g^{\{x,y,z\}}_t$ analogously.
$\bar x=\tfrac1L\sum_t x_t$ is the mean. $\tilde p$ denotes the total-pressure curve resampled to a
fixed length and divided by its mean (a unit-scale shape). The \emph{flat-foot instant} is
$t^\star=\arg\min_t \lVert g_t\rVert$ (foot momentarily still), where the accelerometer points along
gravity.

Several features are \emph{step-to-step} quantities: for a per-step scalar $u^{(k)}$ (step index
$k$ within a recording+foot sequence), the rolling 3-step standard deviation is
\begin{equation}
\operatorname{roll3\_std}(u)^{(k)}=\sqrt{\tfrac{1}{3}\textstyle\sum_{j=k-2}^{k}\bigl(u^{(j)}-\bar u^{(k)}\bigr)^2},
\qquad \bar u^{(k)}=\tfrac13\textstyle\sum_{j=k-2}^{k}u^{(j)}.
\label{eq:roll3}
\end{equation}

\section{Foot-tilt variability (the strongest features)}
From a ZUPT-gated orientation estimate: at the flat-foot instant $t^\star$ the foot tilt is read
from the accelerometer on two orthogonal axes,
\begin{equation}
\theta_A=\operatorname{atan2}\!\bigl(a^x_{t^\star},\,\sqrt{(a^y_{t^\star})^2+(a^z_{t^\star})^2}\bigr),
\qquad
\theta_B=\operatorname{atan2}\!\bigl(a^z_{t^\star},\,\sqrt{(a^x_{t^\star})^2+(a^y_{t^\star})^2}\bigr).
\end{equation}
\begin{itemize}
  \item \code{orient_rest_a_roll3_std} \Fval{3109} and \code{orient_rest_b_roll3_std} \Fval{1710}
  --- the \textbf{step-to-step variability} of resting foot tilt, i.e.\
  $\operatorname{roll3\_std}(\theta_A)$ and $\operatorname{roll3\_std}(\theta_B)$ via
  Eq.~\eqref{eq:roll3}. High values mean the foot lands at a different orientation each step --- the
  signature of uneven/compliant ground. The single most class-differentiable features, and they
  generalise across people. (The \emph{absolute} tilt $\theta_A,\theta_B$ is a per-person mount
  offset and is \emph{not} a feature --- only its step-to-step variability is kept.)
  \item \code{orient_range_a} / \code{orient_range_b} --- the range of foot tilt \emph{within a
  single stance}, $\max_t\theta_{A,t}-\min_t\theta_{A,t}$: how much the foot rocks while planted.
\end{itemize}

\section{Accelerometer texture / impact (vibration and hardness)}
Let $x$ be a step-length signal, detrended, with real-FFT power spectrum $P_k=\lvert X_k\rvert^2$ at
frequencies $f_k$ ($k=1,\dots,K$). Define the normalised spectrum $\rho_k=P_k/\sum_j P_j$.
\begin{align}
\text{spectral centroid} &= \frac{\sum_k f_k P_k}{\sum_k P_k}, &
\text{HF fraction} &= \frac{\sum_{k:\,f_k\ge 5\,\text{Hz}} P_k}{\sum_k P_k}, \\[2pt]
\text{spectral entropy} &= \frac{-\sum_k \rho_k\ln\rho_k}{\ln K}, &
\text{jerk ratio} &= \frac{\operatorname{RMS}(\Delta x)}{\operatorname{RMS}(x-\bar x)}
 = \frac{\sqrt{\tfrac1{L-1}\sum_t(x_t-x_{t-1})^2}}{\sqrt{\tfrac1L\sum_t(x_t-\bar x)^2}}.
\end{align}
\begin{itemize}
  \item \code{accel_jerk_ratio} \Fval{2311} --- jerk ratio of the vertical axis $a^y$
  (amplitude-invariant); sharp/jerky strikes on hard ground, damped on compliant.
  \item \code{accel_hf_frac} \Fval{819} --- HF fraction of $\lVert a_t\rVert$; higher on firm surfaces.
  \item \code{accel_spec_centroid} \Fval{693} --- spectral centroid of $\lVert a_t\rVert$; shifts up
  with harder impacts.
  \item \code{accel_spec_entropy}  --- spectral entropy of $\lVert a_t\rVert$ (tonal vs
  broadband texture).
  \item \code{accel_zcr} --- zero-crossing rate of $a^y$,
  $\ \dfrac{1}{L-1}\sum_{t}\mathbf 1\!\left[\operatorname{sign}(a^y_t-\bar a^y)\neq
  \operatorname{sign}(a^y_{t-1}-\bar a^y)\right]$.
  \item \code{accel_mag_strike} --- peak $\lVert a_t\rVert$ in the impact region
  $\max_{t\le L/3}\lVert a_t\rVert$.
\end{itemize}

\section{Gyroscope motion (foot rotation \& stability)}
\begin{itemize}
  \item \code{gyro_x_std} \Fval{1507} --- standard deviation of the $x$ angular rate over the step,
  $\ \sigma_{g^x}=\sqrt{\tfrac1L\sum_t (g^x_t-\bar g^x)^2}$: side-to-side foot rotation/wobble.
  \item \code{gyro_x_strike} / \code{gyro_y_strike} / \code{gyro_z_strike} --- angular rate at the
  strike instant (rotational ``kick'' at contact).
  \item \code{gyro_y_std} / \code{gyro_z_std} --- angular-rate std on the other two axes.
  \item \code{gyro_hf_frac}  --- HF fraction (same formula) of $\lVert g_t\rVert$.
  \item \code{gyro_mag_strike} --- $\lVert g_t\rVert$ at the strike instant.
\end{itemize}

\section{Stance timing \& pressure-curve shape}
On the unit-scale total curve $\tilde p$ (session-normalised shape):
\begin{itemize}
  \item \code{stance_duration} \Fval{1352} --- loaded time per step,
  $\#\{t: p_t>\tfrac12\bar p\}$; compliant ground softens/lengthens contact.
  \item \code{loading_rate} / \code{loading_rate_norm} --- peak rise rate $\max_t \Delta p_t$ (raw
  and $/$scale); hard loads faster.
  \item \code{rise_slope} / \code{fall_slope} --- pressure rise and fall slopes.
  \item \code{slope_asym} --- asymmetry between loading and unloading slopes.
  \item \code{time_to_peak_frac} --- $\arg\max_t p_t / L$ (timing of peak loading).
  \item \code{gu_dip_depth} --- mid-stance dip depth of $\tilde p$,
  $\ 1-\dfrac{\min_{t\in\text{mid}}\tilde p_t}{\max_t \tilde p_t}$; deeper on grass.
  \item \code{gu_load_first_frac} --- $\sum_{t\le L/2}\tilde p_t\big/\sum_t \tilde p_t$.
  \item \code{gu_push_peak} --- height of the push-off (second) peak of $\tilde p$.
  \item \code{gu_load_slope} / \code{gu_unload_slope} --- $\max_t \Delta\tilde p_t$ (loading) and
  $\min_t \Delta\tilde p_t$ (unloading).
\end{itemize}

\section{Weight distribution (heel vs forefoot)}
\begin{itemize}
  \item \code{heel_fore_ratio} \Fval{665} --- $\bar h/\bar f$ (front-to-back load).
  \item \code{heel_impulse_norm} \Fval{695} --- $\sum_t h_t/\text{scale}$ (time-under-load at heel).
  \item \code{heel_skew} / \code{heel_kurt} --- skewness and kurtosis of the heel series,
  \[
  \text{skew}=\frac{\tfrac1L\sum_t (h_t-\bar h)^3}{\sigma_h^{3}},\qquad
  \text{kurt}=\frac{\tfrac1L\sum_t (h_t-\bar h)^4}{\sigma_h^{4}} .
  \]
  \item \code{heel_ripple_frac} --- high-frequency ``ripple'' fraction of the heel signal.
  \item \code{midstance_diff} --- $h_{t^\star}-f_{t^\star}$ at the flat-foot instant.
  \item \code{gu_fore_early_frac} --- forefoot loading share in early stance,
  $\sum_{t\le t_0}\tilde f_t/\sum_t\tilde f_t$.
  \item \code{gu_fore_heel_toff} --- $(\arg\max_t \tilde f_t-\arg\max_t \tilde h_t)/L$.
\end{itemize}

\section{Pressure-map distribution \& dynamics (across the 12 pads)}
Treat the 12 pads as a spatial map. Let $v_i$ be a per-pad load (e.g.\ mean over stance, or the
flat-foot value) and $q_i=v_i/\sum_j v_j$ its normalised share.
\begin{align}
\text{spatial entropy} &= \frac{-\sum_{i=1}^{12} q_i\ln q_i}{\ln 12}, &
\text{participation} &= \frac{1}{12}\cdot\frac{1}{\sum_{i=1}^{12} q_i^{2}} .
\end{align}
\begin{itemize}
  \item \code{gu_spatial_entropy_stance} \Fval{976} and \code{gu_spatial_entropy_ff} --- normalised
  Shannon entropy of the pad distribution (whole stance / at flat-foot). High $=$ load spread
  evenly; low $=$ concentrated. Compliant ground conforms and spreads load.
  \item \code{gu_participation_ff} --- normalised inverse participation ratio (effective fraction of
  pads bearing load at flat-foot).
  \item \code{gu_active_sensors} --- $\#\{i: v_i>0.1\max_j v_j\}$ at flat-foot (contact breadth).
  \item \code{gu_cop_wander} --- mean frame-to-frame change of the normalised map,
  $\ \tfrac1{L-1}\sum_t\sum_i \lvert Q_{i,t}-Q_{i,t-1}\rvert$ with $Q_{\cdot,t}$ the per-frame
  normalised map: how much load shifts around the sole.
  \item \code{gu_press_jitter} --- mean per-pad temporal jitter,
  $\operatorname{mean}_i \operatorname{std}_t(\Delta^2 s_{i,t})/\text{scale}$.
  \item \code{stance_duration_roll3_std} / \code{loading_rate_roll3_std} /
  \code{impact_mag_roll3_std} --- step-to-step variability (Eq.~\eqref{eq:roll3}) of stance
  duration, loading rate, and impact.
\end{itemize}

\section*{Dropped (computed but excluded from the model)}
Removed because cross-user permutation testing showed they act as a \textbf{person fingerprint} and
\emph{hurt} generalisation to a new person:
\begin{itemize}
  \item \textbf{Per-sensor pressure ratios} --- \code{pressure_01_mean_ratio} $\dots$
  \code{pressure_12_mean_ratio}, plus \code{heel_mean_ratio}, \code{fore_mean_ratio},
  \code{heel_pmax_ratio}, \code{fore_pmax_ratio}.
  \item \textbf{Raw magnitudes} --- \code{total_p_mean}, \code{total_pmax}, \code{impact_mag},
  \code{impact_mag_norm}, \code{heel_impulse} (scale with body weight / donning tightness).
\end{itemize}

\section*{How to read the model, in one line}
Soft-vs-hard is driven by \textbf{step-to-step foot-tilt variability}, \textbf{accelerometer
jerk/texture}, \textbf{gyro rotation variability}, \textbf{stance timing}, and \textbf{pressure-map
spread} --- all direct measures of surface \emph{compliance}. Because compliance is exactly what a
pressure/IMU insole senses well, this two-class model is the reliable, deployable product.

\end{document}
